%! Inverse Functions — Unit 1, Lesson 1.6 — Slides (Beamer)
\documentclass[aspectratio=169, 11pt]{beamer}
\usepackage{precalculus-beamer}   % provides \plumheader{} and \sectionlabel[color]{}

\begin{document}

% TITLE SLIDE
{
\setbeamercolor{background canvas}{bg=plum}
\begin{frame}[plain]
  \vspace{2.8cm}\hspace{0.55cm}%
  \begin{minipage}{0.9\textwidth}
    {\color{goldacc}\bfseries\small LESSON 1.6}\\[0.5cm]
    {\color{white}\bfseries\LARGE Inverse Functions}\\[1.2cm]
    {\color{lilacmid}\small Precalculus~$\cdot$~Unit 1: Functions and Change}
  \end{minipage}
\end{frame}
}

% HOOK SLIDE
\begin{frame}[t, plain]
  \plumheader{Hook: Two \$14 Parking Receipts}
  \sectionlabel[goldacc]{A Receipt Only Records the Answer}

  Each receipt says what was paid. Neither says how long the car was parked ---
  and that is the only thing you want to know.

  \smallskip
  \begin{columns}[T, onlytextwidth]
    \column{0.46\textwidth}
      \begin{block}{Garage A: \(P(x) = 3x + 2\)}
        A \$14 receipt. How many hours?

        \smallskip
        \(3x + 2 = 14 \;\Rightarrow\; x = \mathbf{4}\) --- and no other number
        of hours gives \$14.
      \end{block}

    \column{0.50\textwidth}
      \begin{block}{Garage B: the posted price list}
        \renewcommand{\arraystretch}{1.2}
        \begin{tabular}{@{}lccccc@{}}
          hours   & 1 & 2 & 3 & 4 & 5 \\
          charges & \$8 & \textbf{\$14} & \textbf{\$14} & \$20 & \$26 \\
        \end{tabular}

        \smallskip
        A \$14 receipt. How many hours? \textbf{The receipt does not say.}
      \end{block}
  \end{columns}

  \medskip
  \begin{center}
    One answer gives the question back. The other does not.
  \end{center}
\end{frame}

% SECTION 1 — GIVING BACK THE QUESTION
\begin{frame}[t, plain]
  \plumheader{1. Giving Back the Question}
  \sectionlabel[redacc]{The Whole Lesson in One Picture}

  \begin{center}
  \begin{tikzpicture}[x=1.2cm, y=1.2cm]
    \draw[thick, plum, fill=lilac, rounded corners=2pt]
      (-3.4,-0.55) rectangle (-1.6,0.55);
    \node at (-2.5,0) {\large \(a\)};
    \draw[thick, plum, fill=goldbg, rounded corners=2pt]
      (1.6,-0.55) rectangle (3.4,0.55);
    \node at (2.5,0) {\large \(b\)};
    \draw[->, very thick, plum] (-1.5,0.32) -- (1.5,0.32);
    \node[above, font=\small] at (0,0.36) {\(f\)};
    \draw[->, very thick, goldacc] (1.5,-0.32) -- (-1.5,-0.32);
    \node[below, font=\small] at (0,-0.36) {\(f^{-1}\)};
    \node[below, font=\small\itshape, text=slate] at (-2.5,-0.62)
      {the question};
    \node[below, font=\small\itshape, text=slate] at (2.5,-0.62)
      {the answer};
  \end{tikzpicture}
  \end{center}

  \medskip
  \begin{center}
    \textbf{The inverse gives back the question.}
    If \(f(a) = b\), then \(f^{-1}(b) = a\).
  \end{center}

  \medskip
  \begin{block}{The signature error of this lesson}
    \(f^{-1}(x) \ne \dfrac{1}{f(x)}\). \quad
    \(P^{-1}(14) = \mathbf{4}\), \quad but \quad \(\dfrac{1}{P(14)} =
    \dfrac{1}{44}\).

    \smallskip
    The reciprocal does arithmetic \textit{to the answer}. The inverse goes back
    to the \textit{question}.
  \end{block}
\end{frame}

% SECTION 2 — THE SWAP
\begin{frame}[t, plain]
  \plumheader{2. The Answers Become the Questions}
  \sectionlabel[goldacc]{Trade the Rows --- Compute Nothing}

  \begin{columns}[T, onlytextwidth]
    \column{0.46\textwidth}
      \begin{center}
      \renewcommand{\arraystretch}{1.25}
      \begin{tabular}{|l|c|c|c|c|c|}
      \hline
      \(x\) (h) & 1 & 2 & 3 & 4 & 5 \\
      \hline
      \(P(x)\) (\$) & 5 & 8 & 11 & 14 & 17 \\
      \hline
      \end{tabular}

      \smallskip
      \begin{tabular}{|l|c|c|c|c|c|}
      \hline
      \(x\) (\$) & 5 & 8 & 11 & 14 & 17 \\
      \hline
      \(P^{-1}(x)\) (h) & 1 & 2 & 3 & 4 & 5 \\
      \hline
      \end{tabular}
      \end{center}

      \smallskip
      The questions you may ask \(P^{-1}\) are exactly the answers \(P\) gave
      --- so \textbf{domain and range swap}.

    \column{0.50\textwidth}
      \begin{center}
      \begin{tikzpicture}[x=0.33cm, y=0.33cm]
        \draw[linegray, very thin, step=1] (-5,-5) grid (5,5);
        \draw[->, thick] (-5.4,0) -- (5.7,0) node[below right] {\small \(x\)};
        \draw[->, thick] (0,-5.4) -- (0,5.7) node[above] {\small \(y\)};
        \draw[dashed, thick, slate] (-5,-5) -- (5,5);
        \node[font=\tiny, slate, above left] at (4.6,4.6) {\(y = x\)};
        \draw[very thick, plum] (-3,-5) -- (2,5);
        \draw[very thick, goldacc] (-5,-3) -- (5,2);
        \fill[plum] (0,1) circle (2.2pt);
        \fill[goldacc] (1,0) circle (2.2pt);
        \fill[plum] (2,5) circle (2.2pt);
        \fill[goldacc] (5,2) circle (2.2pt);
      \end{tikzpicture}
      \end{center}

      \smallskip
      Questions live on the \(x\)-axis, answers on the \(y\)-axis. Trading them
      sends \((a,b)\) to \((b,a)\) --- a reflection over \(y = x\).
  \end{columns}
\end{frame}

% SECTION 3 — NO INVERSE
\begin{frame}[t, plain]
  \plumheader{3. When the Answer Gives Back Nothing}
  \sectionlabel[redacc]{Back to the Second Receipt}

  \begin{columns}[T, onlytextwidth]
    \column{0.48\textwidth}
      \begin{block}{Garage B has no inverse}
        \(B^{-1}(14)\) would have to be \(2\) \textbf{and} \(3\) --- and a
        function gets one answer per question.

        \smallskip
        Nothing is wrong with \(B\). The trouble appears only when you run it
        backwards.
      \end{block}

      \medskip
      \begin{block}{One-to-one}
        No two different questions share an answer. \textbf{A function has an
        inverse exactly when it is one-to-one.}
      \end{block}

    \column{0.48\textwidth}
      \begin{block}{\(f(x) = x^2\)}
        \renewcommand{\arraystretch}{1.2}
        \begin{tabular}{@{}lccccc@{}}
          \(x\)   & \(-2\) & \(-1\) & 0 & 1 & 2 \\
          \(x^2\) & \textbf{4} & 1 & 0 & 1 & \textbf{4} \\
        \end{tabular}

        \smallskip
        The answer 4 came from two questions.
      \end{block}

      \medskip
      \begin{block}{The repair: throw questions away}
        Keep only \(x \ge 0\). Then \(f^{-1}(x) = \sqrt{x}\).

        \smallskip
        ``\(x^2\) has no inverse'' is \textbf{false}. Say the domain.
      \end{block}
  \end{columns}
\end{frame}

% SECTION 4 — SWAP AND SOLVE
\begin{frame}[t, plain]
  \plumheader{4. Swap and Solve}
  \sectionlabel[goldacc]{Every Answer at Once}

  A table only gives back the questions already printed on it.

  \medskip
  \begin{columns}[T, onlytextwidth]
    \column{0.44\textwidth}
      \begin{block}{The four steps}
        \begin{enumerate}\setlength{\itemsep}{2pt}
          \item Write \(y = f(x)\).
          \item \textbf{Swap} \(x\) and \(y\).
          \item \textbf{Solve} for \(y\).
          \item Rename it \(f^{-1}(x)\).
        \end{enumerate}

        \smallskip
        Step 2 is the whole inversion: the old \textit{answer} becomes the new
        \textit{question}.
      \end{block}

    \column{0.52\textwidth}
      \(P(x) = 3x + 2\):
      \begin{align*}
        y &= 3x + 2 \\
        x &= 3y + 2 \quad \text{\small(swap)} \\
        \frac{x-2}{3} &= y
      \end{align*}

      \smallskip
      So \(P^{-1}(x) = \dfrac{x-2}{3}\), and \(P^{-1}(14) = \mathbf{4}\) ---
      the receipt, answered.
  \end{columns}

  \medskip
  \begin{center}
    On the warm-up you solved \(3x + 2 = 14\) and got \textit{one} question
    back. This gets them all.
  \end{center}
\end{frame}

% SECTION 5 — THE CHECK
\begin{frame}[t, plain]
  \plumheader{5. The Check: Hand the Answer Back}
  \sectionlabel[redacc]{Where Lesson 1.5 Pays Off}

  In \(f\big(f^{-1}(x)\big)\) the \textbf{handoff} is whatever \(f^{-1}\)
  produced. If \(f^{-1}\) really gives back questions, the far end is where you
  started.

  \medskip
  \begin{columns}[T, onlytextwidth]
    \column{0.52\textwidth}
      \begin{align*}
        P\big(P^{-1}(x)\big) &= 3\!\left(\frac{x-2}{3}\right) + 2 = x \\[4pt]
        P^{-1}\big(P(x)\big) &= \frac{(3x+2) - 2}{3} = x
      \end{align*}

    \column{0.44\textwidth}
      \begin{block}{Both directions, every time}
        \(f\big(f^{-1}(x)\big) = f^{-1}\big(f(x)\big) = x = I(x)\)

        \smallskip
        The \textbf{identity function} from Lesson 1.5, right on schedule.
      \end{block}
  \end{columns}

  \medskip
  \begin{center}
    Warm-up item 2 was this check: \(f(x) = 3x+2\) and \(g(x) =
    \dfrac{x-2}{3}\) gave back 11 and 3. That \(g\) was \(P^{-1}\) all along.
  \end{center}
\end{frame}

% CLASS PRACTICE
\begin{frame}[t, plain]
  \plumheader{Class Practice}
  \sectionlabel[redacc]{Name the Question, Not the Reciprocal}

  \begin{enumerate}\setlength{\itemsep}{10pt}
    \item A function \(g\) satisfies \(g(-4) = 7\). Find \(g^{-1}(7)\), and say
          in one sentence why it is not \(\tfrac{1}{7}\).

    \item Find \(f^{-1}(x)\) for \(f(x) = 5x - 3\), then check it by computing
          \(f^{-1}\big(f(x)\big)\).

    \item A table gives \(h(1) = 5\), \(h(2) = 8\), \(h(3) = 5\), \(h(4) = 12\).
          Explain why \(h\) has no inverse.
  \end{enumerate}
\end{frame}

\end{document}
